% sections/instructions.tex -- grading legend and general notes

\section*{How to Read This Solution Set}
\addcontentsline{toc}{section}{How to Read This Solution Set}

Each problem below is restated in a shaded box, followed by a numbered worked
solution. Small notes in the right margin, like the one beside this
paragraph, reproduce the marking rubric used by the grading team: they show
how many of the problem's points are attached to a given step and what a
grader is checking for.

\rubric{--}{Margin notes follow this format throughout: point value, then
the grading criterion in one line.}

\begin{center}
\begin{tabular}{@{}p{2.1cm}p{9.2cm}@{}}
\toprule
\textbf{Symbol} & \textbf{Meaning} \\
\midrule
\textcolor{inkblue}{\faPen} & Marking-rubric note (margin) --- points and criterion for the adjacent step \\
\textcolor{correctgreen}{\faCheck} & Final answer, boxed for quick reference \\
\textcolor{accentgold}{\faExclamation} & Common mistake flagged at that step \\
\bottomrule
\end{tabular}
\end{center}

\vspace{0.4cm}
\noindent\textbf{Partial credit policy.} Marks for a given step are awarded
independently of later steps: an arithmetic slip in Step 2 does not forfeit
the marks earned for correct setup in Step 1 or correct method in Step 3,
provided the error is carried forward consistently. Full derivations are
required --- a correct final answer with no supporting work earns at most
half of the marks allotted to the problem.
